\section{Quaternion Algebra}
\label{sec:quaternions}

% ----------------------------------------------------------------
% 7 subsections: motivation, definition, rotation proof,
% geometric interpretation, rotation matrix, integration, advantages.
% ----------------------------------------------------------------

% ============================================================
%  3.1  Motivation: rotations as numbers
% ============================================================
\subsection{Motivation: Rotations as Numbers}
\label{sec:quat:motivation}

Before introducing quaternions, it is instructive to recall how
\emph{complex numbers} encode rotations in two dimensions.  A point
$(x,y)\in\mathbb{R}^2$ can be identified with the complex number
$z = x + iy$.  Multiplying $z$ by the unit complex number
$e^{i\theta} = \cos\theta + i\sin\theta$ yields

\begin{equation}
  z' \;=\; e^{i\theta}\,z
     \;=\; (x\cos\theta - y\sin\theta)
          + i\,(x\sin\theta + y\cos\theta),
  \label{eq:complex_rotation}
\end{equation}

\noindent which is exactly the rotation of $(x,y)$ by an angle $\theta$.
For example, rotating the point $(1,0)$ by $90^\circ$ amounts to multiplying
$z=1$ by $e^{i\pi/2} = i$, which gives $z'= i$, corresponding to the
point $(0,1)$---precisely the expected result.

A natural question arises: \emph{does an analogous number system exist
for rotations in $\mathbb{R}^3$?}  One might expect that three imaginary
units suffice, but William Rowan Hamilton showed in 1843 that a consistent
non-commutative algebra requires \textbf{four} components---one real and
three imaginary---giving rise to the \emph{quaternions}
$\mathbb{H} = \{q_w + q_x\mathbf{i} + q_y\mathbf{j} + q_z\mathbf{k}\}$.
Just as the unit complex numbers $e^{i\theta}$ parametrise $\mathrm{SO}(2)$,
the unit quaternions provide a double cover of $\mathrm{SO}(3)$: every
rotation in three-dimensional space can be expressed as a quaternion
conjugation, as we shall demonstrate in the following subsections.

% ============================================================
%  3.2  Definition and operations  (existing content, preserved)
% ============================================================
\subsection{Definition and Key Operations}
\label{sec:quat:def}

A quaternion $\mathbf{q}$ is a four-dimensional hypercomplex number:

\begin{equation}
  \mathbf{q} = q_w + q_x\,\mathbf{i} + q_y\,\mathbf{j} + q_z\,\mathbf{k}
             = \begin{bmatrix} q_w \\ q_x \\ q_y \\ q_z \end{bmatrix}
  \label{eq:quat_def}
\end{equation}

\noindent where $q_w$ is the \emph{scalar} (or real) part and
$(q_x, q_y, q_z)$ is the \emph{vector} (or imaginary) part. The basis
elements satisfy the fundamental relations $\mathbf{i}^2 = \mathbf{j}^2 =
\mathbf{k}^2 = \mathbf{i}\mathbf{j}\mathbf{k} = -1$.

\paragraph{Norm.}
$\|\mathbf{q}\| = \sqrt{q_w^2 + q_x^2 + q_y^2 + q_z^2}$.
A \emph{unit quaternion} satisfies $\|\mathbf{q}\| = 1$.

\paragraph{Conjugate.}
$\mathbf{q}^* = q_w - q_x\,\mathbf{i} - q_y\,\mathbf{j} - q_z\,\mathbf{k}$.
For unit quaternions, $\mathbf{q}^* = \mathbf{q}^{-1}$.

\paragraph{Hamilton product.}
Given $\mathbf{p} = (p_w, p_x, p_y, p_z)$ and
$\mathbf{q} = (q_w, q_x, q_y, q_z)$, their product $\mathbf{r} =
\mathbf{p} \otimes \mathbf{q}$ is:

\begin{equation}
  \mathbf{r} = \begin{bmatrix}
    p_w q_w - p_x q_x - p_y q_y - p_z q_z \\
    p_w q_x + p_x q_w + p_y q_z - p_z q_y \\
    p_w q_y - p_x q_z + p_y q_w + p_z q_x \\
    p_w q_z + p_x q_y - p_y q_x + p_z q_w
  \end{bmatrix}^\top
  \label{eq:hamilton}
\end{equation}

\noindent Note that the Hamilton product is \emph{not} commutative:
$\mathbf{p} \otimes \mathbf{q} \neq \mathbf{q} \otimes \mathbf{p}$ in general.

% ============================================================
%  3.3  Unit quaternion as a rotation in R^3  (rewritten)
% ============================================================
\subsection{The Unit Quaternion as a Rotation in $\mathbb{R}^3$}
\label{sec:quat:rotation}

A unit quaternion $\mathbf{q}$ encodes a rotation by angle~$\theta$ about a
unit axis~$\hat{\mathbf{u}}$:

\begin{equation}
  \mathbf{q} = \cos\frac{\theta}{2}
             + \sin\frac{\theta}{2}\,
               (u_x\,\mathbf{i} + u_y\,\mathbf{j} + u_z\,\mathbf{k})
  \label{eq:axis_angle}
\end{equation}

To rotate a vector $\mathbf{v} \in \mathbb{R}^3$, we embed it as a pure
quaternion $\mathbf{v}_q = (0, v_x, v_y, v_z)$ and compute:

\begin{equation}
  \mathbf{v}'_q = \mathbf{q} \otimes \mathbf{v}_q \otimes \mathbf{q}^*
  \label{eq:rotation}
\end{equation}

\noindent The vector part of $\mathbf{v}'_q$ is the rotated vector
$\mathbf{v}' \in \mathbb{R}^3$.  We now prove that this conjugation is
indeed a rotation---that is, an element of $\mathrm{SO}(3)$.

\paragraph{Preservation of norms.}
The norm of a quaternion product satisfies
$\|\mathbf{p}\otimes\mathbf{q}\| = \|\mathbf{p}\|\,\|\mathbf{q}\|$.
Since $\|\mathbf{q}\|=\|\mathbf{q}^*\|=1$, we obtain
\begin{equation}
  \|\mathbf{v}'_q\|
  = \|\mathbf{q}\otimes\mathbf{v}_q\otimes\mathbf{q}^*\|
  = \|\mathbf{q}\|\;\|\mathbf{v}_q\|\;\|\mathbf{q}^*\|
  = \|\mathbf{v}_q\|,
  \label{eq:norm_preservation}
\end{equation}
so the map preserves the length of every vector.

\paragraph{Linearity.}
For any scalars $\alpha,\beta\in\mathbb{R}$ and pure quaternions
$\mathbf{v}_q,\mathbf{w}_q$:
\begin{equation}
  \mathbf{q}\otimes(\alpha\mathbf{v}_q+\beta\mathbf{w}_q)\otimes\mathbf{q}^*
  = \alpha\,(\mathbf{q}\otimes\mathbf{v}_q\otimes\mathbf{q}^*)
  + \beta\,(\mathbf{q}\otimes\mathbf{w}_q\otimes\mathbf{q}^*),
  \label{eq:linearity}
\end{equation}
which follows from the bilinearity of the Hamilton product.  Hence the map
$\mathbf{v}\mapsto\mathbf{q}\otimes\mathbf{v}_q\otimes\mathbf{q}^*$
defines a linear transformation $\mathbb{R}^3\to\mathbb{R}^3$.

\paragraph{Preservation of orientation.}
A norm-preserving linear map is an \emph{orthogonal} transformation, so
its determinant is $\pm 1$.  Because the set of unit quaternions is
connected (it is the 3-sphere $S^3$) and $\mathbf{q}=1$ (the identity
quaternion) produces the identity map with $\det=+1$, and the determinant
is a continuous function of $\mathbf{q}$, we conclude that
$\det=+1$ for \emph{every} unit quaternion.  Therefore the conjugation
in~\eqref{eq:rotation} is a \textbf{proper rotation}, i.e.\ an element
of $\mathrm{SO}(3)$.

% ============================================================
%  3.4  Geometric interpretation: parallel/perpendicular decomposition
% ============================================================
\subsection{Geometric Interpretation}
\label{sec:quat:geometry}

The proof above establishes \emph{that} the conjugation is a rotation.
We now show \emph{how} it acts geometrically, which also explains the
factor $\theta/2$ in the quaternion parametrisation.

Let $\hat{\mathbf{u}}$ be the rotation axis and decompose an arbitrary
vector $\mathbf{v}$ into its parallel and perpendicular components with
respect to $\hat{\mathbf{u}}$:
\begin{equation}
  \mathbf{v} = \mathbf{v}_\parallel + \mathbf{v}_\perp,
  \qquad
  \mathbf{v}_\parallel = (\mathbf{v}\cdot\hat{\mathbf{u}})\,\hat{\mathbf{u}},
  \qquad
  \mathbf{v}_\perp = \mathbf{v} - \mathbf{v}_\parallel.
  \label{eq:decomposition}
\end{equation}

\paragraph{Parallel component is invariant.}
Because $\mathbf{v}_\parallel$ is collinear with $\hat{\mathbf{u}}$, and
$\hat{\mathbf{u}}$ commutes with every quaternion whose vector part is
proportional to $\hat{\mathbf{u}}$ (in particular with
$\mathbf{q}=\cos(\theta/2)+\sin(\theta/2)\,\hat{\mathbf{u}}$), we have:
\begin{equation}
  \mathbf{q}\otimes(\mathbf{v}_\parallel)_q\otimes\mathbf{q}^*
  = (\mathbf{v}_\parallel)_q.
  \label{eq:parallel_invariant}
\end{equation}
In other words, the component along the rotation axis is unchanged.

\paragraph{Perpendicular component is rotated by $\theta$.}
For the perpendicular component, the double-sided multiplication produces:
\begin{equation}
  \mathbf{q}\otimes(\mathbf{v}_\perp)_q\otimes\mathbf{q}^*
  = \cos\theta\;\mathbf{v}_\perp
  + \sin\theta\;(\hat{\mathbf{u}}\times\mathbf{v}_\perp),
  \label{eq:perp_rotation}
\end{equation}
which is Rodrigues' rotation formula restricted to vectors orthogonal
to the axis.  The key insight is that \emph{each} multiplication by
$\mathbf{q}$ or $\mathbf{q}^*$ contributes a half-angle rotation, so the
two multiplications together produce the full angle~$\theta$.  This is
precisely why the quaternion uses $\theta/2$ rather than $\theta$: the
\textbf{double-sided conjugation doubles the angle}.

Combining both components, we recover Rodrigues' formula for an arbitrary
vector:
\begin{equation}
  \mathbf{v}' = \cos\theta\;\mathbf{v}
              + (1-\cos\theta)\,(\mathbf{v}\cdot\hat{\mathbf{u}})\,\hat{\mathbf{u}}
              + \sin\theta\;(\hat{\mathbf{u}}\times\mathbf{v}).
  \label{eq:rodrigues}
\end{equation}

% ============================================================
%  3.5  Equivalence with the rotation matrix R(q)
% ============================================================
\subsection{Equivalence with the Rotation Matrix}
\label{sec:quat:matrix}

By expanding the conjugation
$\mathbf{q}\otimes\mathbf{v}_q\otimes\mathbf{q}^*$ component by component
using the Hamilton product~\eqref{eq:hamilton} and grouping terms in
$v_x$, $v_y$, $v_z$, one obtains the equivalent matrix form
$\mathbf{v}' = R(\mathbf{q})\,\mathbf{v}$, where:

\begin{equation}
  R(\mathbf{q}) = \begin{bmatrix}
    1 - 2(q_y^2+q_z^2)  & 2(q_x q_y - q_w q_z)  & 2(q_x q_z + q_w q_y) \\[4pt]
    2(q_x q_y + q_w q_z) & 1 - 2(q_x^2+q_z^2)    & 2(q_y q_z - q_w q_x) \\[4pt]
    2(q_x q_z - q_w q_y) & 2(q_y q_z + q_w q_x)   & 1 - 2(q_x^2+q_y^2)
  \end{bmatrix}
  \label{eq:rotation_matrix}
\end{equation}

\noindent This matrix is used whenever an explicit $3\times 3$ rotation is
needed---for instance, to project accelerometer readings into the world
frame.

\paragraph{Verification.}
One can verify that $R(\mathbf{q})$ is orthogonal and proper:
\begin{enumerate}
  \item \textbf{Orthogonality}: $R(\mathbf{q})^\top R(\mathbf{q}) = I$,
    which follows from the unit-norm constraint $q_w^2+q_x^2+q_y^2+q_z^2=1$.
  \item \textbf{Unit determinant}: $\det R(\mathbf{q}) = +1$, confirming
    that $R(\mathbf{q})\in\mathrm{SO}(3)$.
\end{enumerate}

\noindent Note that $\mathbf{q}$ and $-\mathbf{q}$ yield the same
rotation matrix $R(\mathbf{q}) = R(-\mathbf{q})$, reflecting the
\emph{double cover} of $\mathrm{SO}(3)$ by the unit quaternions.

% ============================================================
%  3.6  Integration from angular velocity  (existing content, preserved)
% ============================================================
\subsection{Quaternion Integration from Angular Velocity}
\label{sec:quat:integration}

Given the current orientation quaternion $\mathbf{q}(t)$ and a gyroscope
reading $\boldsymbol{\omega}_b = (\omega_x, \omega_y, \omega_z)$, the
orientation evolves according to the differential equation:

\begin{equation}
  \dot{\mathbf{q}}(t) = \frac{1}{2}\,\mathbf{q}(t) \otimes
  \boldsymbol{\omega}_q
  \label{eq:quat_diff}
\end{equation}

\noindent where $\boldsymbol{\omega}_q = (0, \omega_x, \omega_y, \omega_z)$ is
the angular velocity expressed as a pure quaternion. Using a first-order
discretisation with time step~$\Delta t$:

\begin{equation}
  \mathbf{q}(t + \Delta t) = \mathbf{q}(t) + \frac{\Delta t}{2}\,
  \mathbf{q}(t) \otimes \boldsymbol{\omega}_q
  \label{eq:quat_euler}
\end{equation}

\noindent After each update, $\mathbf{q}$ must be renormalised to remain on the
unit sphere: $\mathbf{q} \leftarrow \mathbf{q} / \|\mathbf{q}\|$.

% ============================================================
%  3.7  Advantages over Euler angles and matrices
% ============================================================
\subsection{Advantages over Euler Angles and Rotation Matrices}
\label{sec:quat:advantages}

Table~\ref{tab:rotation_comparison} summarises the key properties of the
three most common rotation representations.  Quaternions offer the best
trade-off between compactness, numerical stability, and computational cost
for real-time sensor fusion applications.

\begin{table}[t]
  \centering
  \caption{Comparison of rotation representations.}
  \label{tab:rotation_comparison}
  \small
  \begin{tabular}{@{}lccc@{}}
    \toprule
    \textbf{Property} & \textbf{Euler} & \textbf{Matrix} & \textbf{Quaternion} \\
    \midrule
    Parameters             & 3     & 9     & 4 \\
    Gimbal lock            & Yes   & No    & No \\
    Singularity-free       & No    & Yes   & Yes \\
    Composition cost       & High  & $O(27)$ & $O(16)$ \\
    Renormalisation        & ---   & $O(n^3)$ & $O(n)$ \\
    Interpolation (SLERP)  & Poor  & Poor  & Native \\
    \bottomrule
  \end{tabular}
\end{table}

Euler angles, while intuitive, suffer from \emph{gimbal lock}---a
singularity that occurs when two rotation axes become aligned, causing a
loss of one degree of freedom.  Rotation matrices are singularity-free
but use nine parameters (with six orthogonality constraints), making
renormalisation expensive.  Quaternions use only four parameters (with
one constraint $\|\mathbf{q}\|=1$), avoid gimbal lock entirely, compose
efficiently via the Hamilton product, and support smooth interpolation
through spherical linear interpolation (SLERP).  These properties make
quaternions the natural choice for the sensor fusion pipeline developed
in this work.
